\begin{abstract}
\textbf{Background:} Opportunistic bacterial pathogens colonizing life support systems and silicone water conduits aboard spaceflight platforms pose severe health risks to crew members during long-duration interplanetary missions. In NASA Open Science Data Repository (OSDR) study \textbf{OSD-528} (Clary et al., \textit{Front. Space Technol.} 2022), the fish pathogen and surrogate for \textit{Mycobacterium tuberculosis}, \textit{Mycobacterium marinum} 1218R, was cultured on hydrophobic polydimethylsiloxane (PDMS) silicone membranes under simulated microgravity using a custom low-cost 3D clinostat and a Random Positioning Machine (RPM 2.0) alongside static 1g ground controls. However, systems-level classification and biological pathway prioritization have been constrained by the small-sample bottleneck ($N=9$) typical of space biology.

\textbf{Methods:} We established a FAIR-compliant (Findable, Accessible, Interoperable, Reusable) systems biology framework that couples Weighted Gene Co-expression Network Analysis (WGCNA) with the newly introduced \textbf{TabPFN} tabular foundation model (\textit{Nature} 2025). WGCNA collapsed 5,424 \textit{M. marinum} genes into 5 discrete topological modules ($\beta=6$), extracting Module Eigengenes (MEs) and intramodular hub centralities ($k_{\text{within}}$). To validate the pipeline, an empirically parameterized simulation benchmark derived from published OSD-528 phenotypes was coupled with TabPFN in-context learning to perform leave-one-out cross-validation (LOOCV) modality classification, benchmarked against Random Forest, alongside permutation feature importance and multi-scale Gene Ontology enrichment.

\textbf{Results:} TabPFN achieved 100.0\% LOOCV classification accuracy across the three simulation modalities (3D Clinostat, RPM 2.0, Static 1g) and 100.0\% binary microgravity detection, significantly outperforming classical Random Forest (44.4\%). Permutation importance identified discrete kinematic markers separating simulator mechanics (isocitrate lyase \textit{icl1}, chaperone \textit{dnaK}, protease \textit{clpP1}, and cytochrome $bd$ oxidase \textit{cydA}). Simultaneously, WGCNA revealed a massive, concordant core microgravity network ($r > 0.97, p < 10^{-30}$) shared between 3D Clinostat and RPM 2.0, governed by glycopeptidolipid biofilm synthetases (\textit{mps1}, \textit{mps2}, \textit{fadD28}), mycolic acid FAS-II cell envelope elongases (\textit{kasA}, \textit{kasB}, \textit{acpM}, \textit{inhA}, \textit{fbpA}), Type VII ESX-1 virulence secretion machinery (\textit{esxA}, \textit{esxB}, \textit{eccA1--eccE1}), and the DosR hypoxia/dormancy regulon (\textit{dosR}, \textit{hspX}).

\textbf{Conclusions:} This study demonstrates that low-cost 3D clinorotation and RPM 2.0 simulate virtually identical core mycobacterial cell envelope and virulence adaptations, validating the utility of accessible ground hardware for spaceflight risk mitigation. By combining WGCNA dimensionality reduction with tabular foundation AI, we establish a generalizable blueprint for small-sample spaceflight meta-analyses, released with complete RO-Crate v1.1 and Zenodo metadata.
\end{abstract}
